% Polygonal Jordan theorem.
\section{Jordan's theorem for polygons}

A \emph{polygon} $P$ is a closed path $v_1v_2\cdots v_nv_1$ made of $n\ge3$ segments (the \emph{edges} $e_i=[v_i,v_{i+1}]$, indices modulo $n$). It is \emph{simple} if two edges meet only when
they are consecutive, and then only at their common vertex; consecutive edges are not collinear. We identify $P$ with the union of its edges, a compact subset of $\R^2$.

\begin{theorem}[Jordan's theorem for polygons]\label{thm:jordanpoly}
If $P$ is a simple polygon, then $\R^2\setminus P$ has exactly two connected components. One is bounded and the other is not, and $P$ is the boundary of each.
\end{theorem}

\begin{figure}[ht]\centering\fig{ch10-polygon}
\caption{The parity of the number of crossings of a ray with $P$ tells inside from outside: $3$ for the point $x$, $2$ for the point $y$.}\label{fig:polygon}
\end{figure}

The proof has two halves. Counting crossings of a ray (Figure~\ref{fig:polygon}) shows that there are \emph{at least} two components; walking along both sides of $P$ shows that there are \emph{at most} two.
Here ``connected'' means path-connected, and the components of an open subset of the plane are open.

\subsection*{At least two components: a parity}

Choose a unit vector $d$ that is parallel to no edge and to no line through two vertices; only finitely many directions are excluded. For a vertex $v$, call the line $v+\R d$ a \emph{bad line}; there are $n$ of them,
and each contains exactly one vertex. For $x\notin P$ off the bad lines, the ray $x+\R_{\ge0}d$ meets no vertex; it meets each edge in at most one point (it is not parallel to the edge), and there it crosses the
edge from one side to the other. Let $c(x)$ be the number of edges it meets, and $\pi(x)=c(x)\bmod2$.

\begin{lemma}\label{lem:parity}
Every $x\in\R^2\setminus P$ has a disc neighbourhood $D$, disjoint from $P$, on which $\pi$ is constant (on the points of $D$ off the bad lines).
\end{lemma}

\begin{proof}
Let $x\notin P$. The ray from $x$ meets at most one vertex. Let $\rho>0$ be smaller than the distance from $x$ to $P$, and from the ray at $x$ to every vertex not on it. Then for $x'$ in the disc $D$ of
radius $\rho$ about $x$, the ray from $x'$ is a translate of the ray from $x$ by less than $\rho$.

If the ray from $x$ meets no vertex, then (taking $\rho$ smaller if needed, by compactness of the edges and transversality) the ray from $x'$ meets exactly the edges met by the ray from $x$, and $\pi(x')=\pi(x)$. If it meets the vertex $v$,
between the edges $e$ and $f$, the same holds for all edges except $e$ and $f$. Points $x'$ of $D$ off the bad line through $v$ lie on one of its two sides, and the ray from $x'$ passes $v$ on that side. An edge at $v$ is met by
the ray from $x'$ exactly when the edge leaves $v$ towards the side of the bad line where $x'$ is... more simply: when $x'$ crosses the bad line, the ray sweeps across $v$, and it stops meeting the edges at $v$ on one side and starts
meeting those on the other side. If $e$ and $f$ lie on the same side of the bad line, the count of these two edges changes from $0$ to $2$ or from $2$ to $0$; if they lie on opposite sides, it stays $1$. In all cases $\pi$ does not change.
\end{proof}

By the lemma, $\pi$ extends to a function on all of $\R^2\setminus P$ that is constant on a neighbourhood of every point: put $\pi(x)=\pi(x')$ for any $x'$ near $x$ off the bad lines. A locally constant function is constant along
any path (the set of times where it equals its initial value is open and closed in $[0,1]$), hence on each component.

Now take a point $m$ in the interior of an edge $e$, off the bad lines, and $\varepsilon>0$ so small that the segment $[m-\varepsilon d,m+\varepsilon d]$ meets $P$ only at $m$ and its endpoints are off the bad lines. The ray from $m-\varepsilon d$ is the
segment $[m-\varepsilon d,m+\varepsilon d)$ followed by the ray from $m+\varepsilon d$, and it meets $P$ in one more point, $m$. So $\pi(m-\varepsilon d)\ne\pi(m+\varepsilon d)$: these two points lie in different components, and there are at least two.

\subsection*{At most two components: the two sides}

Let $r>0$ be smaller than half the length of every edge and than half the distance from each vertex to every edge not containing it (positive, since $P$ is simple). Let $B_v$ be the open disc of radius $r$ about the vertex $v$.
In polar coordinates about $v$, $B_v\cap P$ is two radii, in the directions of the two edges at $v$. They cut $B_v\setminus P$ into two open circular sectors, each path-connected (it is an interval of angles times $(0,r)$).

For an edge $e=[v,w]$, let $e'$ be the part of $e$ at distance at least $\frac r2$ from both $v$ and $w$. It is a compact segment disjoint from all other edges, so it has positive distance $\delta_e$ from them. Let $R_e$ be the open rectangle of points whose orthogonal projection on the line of $e$ lies in the interior of $e'$ and whose distance to that line is less than $h$, with $h<\min(\delta_e,\frac r4)$ and so small that the ends of $R_e$ near $v$ and $w$ stay inside $B_v$ and $B_w$ away from the other edge there. Then $R_e\cap P=e'$, and $R_e\setminus P$ is two open
rectangles $R_e^+$ and $R_e^-$, one on each side of $e$, each convex.

\begin{lemma}\label{lem:sides}
Each of $R_e^\pm$ meets exactly one sector at each end; at each vertex, the two sides of an edge meet different sectors.
\end{lemma}

\begin{proof}
The end of $R_e^+$ near $v$ is a connected subset of $B_v\setminus P$, so it lies in one sector. In polar coordinates about $v$, the points of $R_e^\pm$ near $v$ have angles just above and just below the direction of $e$. The two directions of the edges at $v$ cut the circle of
angles into two arcs, and the direction of $e$ is an endpoint of both; angles just above it lie in one arc and angles just below in the other. So $R_e^+$ and $R_e^-$ meet different sectors.
\end{proof}

Build a graph whose $4n$ nodes are the $2n$ sides $R_e^\pm$ and the $2n$ sectors, with an edge when they intersect. By Lemma~\ref{lem:sides}, each side has degree $2$ (one sector at each end), and each sector, at a vertex where the edges $e$ and $f$ meet, contains exactly one side of $e$ and one side of $f$, so it has degree $2$ too.
A graph in which every node has degree $2$ is a union of disjoint cycles. Start from $R_{e_1}^+$ and walk: sector at $v_2$, a side of $e_2$, sector at $v_3$, and so on around $P$, back to a side of $e_1$ after visiting one side of every edge. Either the walk returns to $R_{e_1}^+$, and the
other $2n$ nodes form a second cycle, or it returns to $R_{e_1}^-$ and continues through the remaining nodes. So there are at most two cycles. The union of the nodes of a cycle is path-connected (each node is, and consecutive nodes intersect), so we get at most two path-connected sets $S_1,S_2\subseteq\R^2\setminus P$ (possibly equal).

The sets $B_v$ and $R_e$ are open and cover $P$, and each of them minus $P$ is contained in $S_1\cup S_2$. Let $U$ be a component of $\R^2\setminus P$. It is open and not the whole plane, and the plane is connected, so $U$ has a boundary point $p$, which lies in $P$ (a point of another
component has a neighbourhood disjoint from $U$). The open set of the cover containing $p$ meets $U$, so $U$ meets $S_1$ or $S_2$, and therefore contains it. Hence there are at most two components.

\begin{proof}[Proof of Theorem~\ref{thm:jordanpoly}]
There are at least two components and at most two, so exactly two, and each contains one of $S_1$, $S_2$; in particular $S_1\ne S_2$, so we are in the case of two cycles, and each component contains one side of every edge and one sector at every vertex. Hence every point of $P$ is a limit
of points of each component, and the boundary of each component, which lies in $P$, is all of $P$. Finally, $P$ lies in a disc $B$; the complement of $B$ is path-connected and disjoint from $P$, so it lies in one component, the unbounded one, and the other component is contained in $B$.
\end{proof}

\begin{remark}[the general theorem]\label{rem:jordangeneral}
For an arbitrary continuous simple closed curve the theorem is true but much harder: one approximates the curve by polygons, and the difficulty is to control the components in the limit. Hales's 2007 formal proof in HOL Light follows this route,
and his article on Jordan's original text explains which steps Jordan took for granted. The polygonal case already shows what readers of Jordan argued about: the treatment of rays through vertices (Lemma~\ref{lem:parity}), which Jordan regarded as obvious.
\end{remark}

Look back at the proof of Lemma~\ref{lem:parity}: it contains a sentence that starts one way and then says ``more simply''. That is how first drafts look, and it is the first thing to fix (Exercise~\ref{ex:rewritelemma}).

\begin{think}
Where exactly was simplicity of $P$ used? Try the proof on a figure eight and see which step fails, and how many components there are.
\end{think}

\begin{beauty}
Jordan's theorem is the archetype of a statement that is obvious and hard. The polygonal proof is honest about why: all the difficulty sits at the vertices, in one parity check and one
picture of two sectors, and once those are written down carefully there is nothing left to believe. Writing well is the art of putting the difficulty where the reader can see it.
\end{beauty}

\begin{exercise}[$\star$]\label{ex:rewritelemma}
Rewrite the second paragraph of the proof of Lemma~\ref{lem:parity} so that it says what it means the first time, with the case analysis at the vertex stated as a separate claim.
\end{exercise}


\section*{Hints}

\begin{hint}{ex:rewritelemma}
Isolate the claim: ``as $x'$ crosses the bad line through $v$, the number of edges at $v$ met by the ray changes by $0$ or $\pm2$''.
\end{hint}

\section*{Solutions}

\begin{solution}{ex:rewritelemma}
\idea State the local claim first, then use it.
\writeup ``Suppose the ray from $x$ meets the vertex $v$, where the edges $e$ and $f$ meet. Every other edge is met by the ray from $x'$ if and only if it is met by the ray from $x$. \emph{Claim:} the number of edges among $e,f$ met by the ray from $x'$ is the same, modulo $2$, on both sides of the bad line through $v$. Indeed, the ray from $x'$ is parallel to the bad line and passes $v$ on the side of $x'$, at small distance; so an edge at $v$ is met by it exactly when the edge lies on the side of $x'$. If $e$ and $f$ lie on the same side, the count is $0$ on one side of the bad line and $2$ on the other; if they lie on opposite sides, it is $1$ on both. Hence $\pi$ is constant on $D$.'' The original sentence tried to give the reason and the conclusion at once; separating the claim from its proof is almost always the fix.
\end{solution}
